\section{Composite Systems and Tensor Products}
While a single isolated qubit provides the fundamental building block of quantum information, the true computational power of the quantum paradigm emerges only when multiple qubits are combined and allowed to interact.

To mathematically describe the collective state space of a multi-partite quantum system, we utilize the \textit{tensor product} (often referred to as the Kronecker product in matrix algebra), denoted by the symbol $\otimes$ \cite{nielsen2010quantum}.

If an isolated physical system $A$ is completely described by a Hilbert space $\mathcal{H}_A$ of dimension $N$, and a second isolated system $B$ is described by a Hilbert space $\mathcal{H}_B$ of dimension $M$, their newly formed composite system $AB$ resides mathematically in the tensor product space $\mathcal{H}_A \otimes \mathcal{H}_B$.
Crucially, the total dimension of this combined space is the multiplicative product $N \times M$, rather than the additive sum $N + M$ typically seen when bolting classical systems together.

For a simple register of two interacting qubits, where each individually resides in $\mathbb{C}^2$, the composite two-qubit state exists within a four-dimensional complex vector space, $\mathbb{C}^2 \otimes \mathbb{C}^2 \cong \mathbb{C}^4$.

The computational basis for this joint system is constructed by taking every possible pairwise combination of the individual basis vectors:
\begin{equation}
    \{\ket{0} \otimes \ket{0}, \ket{0} \otimes \ket{1}, \ket{1} \otimes \ket{0}, \ket{1} \otimes \ket{1}\}
\end{equation}

For notational brevity, it is standard practice to entirely drop the explicit tensor product symbol and condense these basis kets into joint strings, writing them simply as $\ket{00}, \ket{01}, \ket{10},$ and $\ket{11}$.
Consequently, the most general pure state of a two-qubit register, $\ket{\psi}$, is a vast superposition extending across all four of these discrete basis states:
\begin{equation}
    \ket{\psi} = \alpha_{00}\ket{00} + \alpha_{01}\ket{01} + \alpha_{10}\ket{10} + \alpha_{11}\ket{11}
\end{equation}
Just as with a single qubit, the system must definitively exist somewhere upon measurement, forcing the absolute sum of the squared probability amplitudes to equal exactly 100\%: $\sum_{i,j \in \{0,1\}} |\alpha_{ij}|^2 = 1$.

This multiplicative tensor structure reveals a profound physical truth: adding just a single qubit to a register effectively doubles the total number of amplitudes required to fully track the system's state.
For a register consisting of $n$ qubits, the state space explodes to an astonishing $2^n$ complex dimensions.
This exponential scaling mathematically underpins the massive parallelism intrinsic to quantum algorithms, representing a primary source of their theoretical computational supremacy over classical hardware \cite{arora2009computational}.